\documentclass[11pt]{article}
\usepackage[T1]{fontenc}
\usepackage[a4paper,margin=1in]{geometry}
\usepackage{amsmath,amssymb,amsthm}
\usepackage[hidelinks]{hyperref}
\usepackage{microtype}

\newtheorem*{question}{Source Question 2.4 (Kleinberg)}
\newtheorem*{answer}{Answering Theorem 1.3}

\setlength{\parindent}{0pt}
\setlength{\parskip}{0.6em}

\title{Literature Status: Random Regular Graphs\
as Nonlinear Expanders with Respect to Each Other}
\author{Run \texttt{fa\_banach\_001} --- lane 9}
\date{11 August 2026}

\begin{document}
\maketitle

\begin{center}
\fbox{\parbox{0.9\textwidth}{
\textbf{Status: literature already answered.}
Question~2.4 of Mendel--Naor (arXiv:1306.5434) asks whether two
independent random cubic graphs have a dimension-free nonlinear
Poincar\'e constant.  Altschuler--Dodos--Tikhomirov--Tyros
(arXiv:2506.17433, Theorem~1.3) explicitly give a complete affirmative
answer, with arbitrary fixed degrees and all exponents $p\geq1$.
}}
\end{center}

\section{Original question}

Let $\mathcal G_{n,d}$ denote the uniform probability distribution on
simple $d$-regular graphs on $n$ vertices.  For a graph $H$, write $d_H$
for its shortest-path metric, and let $\gamma(G,d_H^2)$ be the least
constant in the corresponding nonlinear Poincar\'e inequality.

In Section~2.1, Question~2.4, on printed page~11 of
\emph{Expanders with respect to Hadamard spaces and random graphs},
Mendel and Naor ask~\cite{source}:

\begin{question}
Does there exist a universal constant $K<\infty$ such that
\[
 \lim_{\min\{m,n\}\to\infty}
 (\mathcal G_{m,3}\!\times\!\mathcal G_{n,3})
 \bigl\{(G,H):\gamma(G,d_H^2)\leq K\bigr\}=1?
\]
\end{question}

Thus the question permits different vertex-set sizes and quantifies over
all maps from the vertices of $G$ into the graph metric of $H$.  The source
proves only the special case of bijections between equal-size graphs in
Proposition~8.1.

\section{Separate later answer}

Altschuler, Dodos, Tikhomirov, and Tyros explicitly restate this as
Problem~1.2 of \emph{Discrete Poincar\'e inequalities and universal
approximators for random graphs}~\cite{answerpaper}.  Their abstract calls
their result a complete affirmative resolution of Kleinberg's problem.
Theorem~1.3, on printed page~3, proves the following stronger statement.

\begin{answer}
For fixed integers $d,\Delta\geq3$ and every $p\geq1$, there are constants
$\Gamma(d,p)<\infty$ and $\tau(d,\Delta)>0$ such that independent
$G\sim\mathcal G_{n,d}$ and $H\sim\mathcal G_{m,\Delta}$ satisfy
\[
 \gamma(G,d_H^p)\leq \Gamma(d,p)
\]
simultaneously for every $p\geq1$, with probability at least
$1-O_{d,\Delta}(\min\{n,m\}^{-\tau})$.
\end{answer}

Taking $d=\Delta=3$ and $p=2$ gives the source statement with
$K=\Gamma(3,2)$, because the displayed error tends to zero as
$\min\{m,n\}\to\infty$.  The match is exact and author-explicit, not an
agent-inferred reformulation.

\section{Scope and search evidence}

The later theorem completely answers Question~2.4 and also strengthens it
to arbitrary fixed degrees and all $p\geq1$.  It does not settle the source
paper's distinct union-stability question for nonlinear spectral gaps in the
remark on printed page~26.

A bounded search through 11 August 2026 checked the run's registry and
solution indexes, the exact wording ``two stochastically independent random
3-regular graphs'', the source title and citations, and recent work on
random-graph nonlinear spectral gaps.  The decisive match is
arXiv:2506.17433v2, whose Problem~1.2 cites Mendel--Naor Question~2.4 and
whose Theorem~1.3 is labelled an affirmative solution.  No new mathematical
result is claimed here.

\begin{thebibliography}{9}
\bibitem{source}
M. Mendel and A. Naor,
\emph{Expanders with respect to Hadamard spaces and random graphs},
Duke Math. J. \textbf{164} (2015), 1471--1548;
arXiv:1306.5434.

\bibitem{answerpaper}
D. J. Altschuler, P. Dodos, K. Tikhomirov, and K. Tyros,
\emph{Discrete Poincar\'e inequalities and universal approximators for
random graphs}, arXiv:2506.17433v2 (2025).
\end{thebibliography}

\end{document}

